% Figure: schematic soft-tissue uniaxial stress-strain curve showing
% the three regions (toe, linear, failure) with the underlying
% collagen-microstructure interpretation.
\begin{figure}[htbp]
\centering
\begin{tikzpicture}[
  axis/.style={-{Stealth}, thick, draw=engblack},
  curve/.style={very thick, engred},
  region/.style={draw=none, opacity=0.15},
  annot/.style={font=\scriptsize, align=center},
]

\draw[axis] (0, 0) -- (9, 0) node[right, font=\small] {strain $\varepsilon$};
\draw[axis] (0, 0) -- (0, 6.5) node[above, font=\small] {stress $\sigma$ (MPa)};

\foreach \x/\xt in {1.5/0.02, 3.0/0.04, 4.5/0.06, 6.0/0.08, 7.5/0.10} {
  \draw (\x, 0) -- (\x, -0.08) node[below, font=\scriptsize] {\xt};
}
\foreach \y/\yt in {1/20, 2/40, 3/60, 4/80, 5/100, 5.5/110} {
  \draw (0, \y) -- (-0.08, \y) node[left, font=\scriptsize] {\yt};
}

% Region shading
\fill[engblue, opacity=0.12] (0, 0) rectangle (2.5, 6);
\fill[englime, opacity=0.12] (2.5, 0) rectangle (7.0, 6);
\fill[engred, opacity=0.12] (7.0, 0) rectangle (8.5, 6);

% Toe region: shallow rise
\draw[curve, smooth, samples=40, domain=0:2.5]
  plot (\x, {0.06 * \x * \x});
% Linear region: nearly straight
\draw[curve, smooth, samples=40, domain=2.5:7.0]
  plot (\x, {0.06*2.5*2.5 + 1.15 * (\x - 2.5)});
% Failure region: yield + drop
\draw[curve, smooth, samples=40, domain=7.0:8.0]
  plot (\x, {0.06*2.5*2.5 + 1.15*(7-2.5) + 0.5*(\x - 7) - 5.0*(\x - 7)^2});

% Mark UTS point
\fill[engblack] (7.5, 5.6) circle (2pt);
\node[font=\scriptsize, anchor=south west] at (7.55, 5.65) {UTS};

% Region labels
\node[annot] at (1.25, 5.7) {Toe\\($\varepsilon < 0.02$)};
\node[annot] at (4.75, 5.7) {Linear\\(structural)};
\node[annot] at (7.6, 5.7) {Failure};

% Microstructure cartoons below x-axis
\begin{scope}[yshift=-2.0cm]
% Toe: wavy fibres
\draw[engred, very thick, smooth] (0.3, 0) sin (0.6, 0.15) cos (0.9, 0)
  sin (1.2, -0.15) cos (1.5, 0) sin (1.8, 0.15) cos (2.1, 0);
\node[annot] at (1.25, -0.5) {wavy collagen};
% Linear: straight fibres
\draw[engred, very thick] (3.0, -0.1) -- (6.5, -0.1);
\draw[engred, very thick] (3.0, 0.1) -- (6.5, 0.1);
\node[annot] at (4.75, -0.5) {straightened, stretched};
% Failure: broken
\draw[engred, very thick] (7.0, 0) -- (7.3, 0);
\draw[engred, very thick] (7.5, 0.15) -- (7.7, 0.15);
\draw[engred, very thick] (7.5, -0.15) -- (7.7, -0.15);
\draw[engred, very thick] (7.9, 0) -- (8.2, 0);
\node[annot] at (7.6, -0.5) {fibres failing};
\end{scope}

\end{tikzpicture}
\caption{Schematic uniaxial stress-strain response of a typical
collagenous soft tissue (tendon, ligament, skin). Three regions: the
toe region ($\varepsilon \lesssim 0.02$) where wavy collagen fibres
straighten under small loads with little stiffness; the linear region
($\varepsilon \approx 0.02$ to $0.08$) where the straightened fibres
load uniformly and the modulus is the structural Young's modulus
(typically $0.5$ to $2\,\text{GPa}$ for tendon); and the failure
region where individual fibres begin to fail and the load decreases
even as strain rises. UTS for human Achilles tendon is in the
$50$ to $100\,\text{MPa}$ range. The toe region is the engineering
reason that small tendon deformations cost the body almost no force,
and is the feature that allows tendons to function as efficient
elastic-energy storage elements during locomotion.}
\label{fig:vol06ch06:soft-stress-strain}
\end{figure}
